\documentclass[11pt]{article}
\usepackage[margin=0.9in]{geometry}
\usepackage{amsmath,amssymb,amsthm,booktabs,mathrsfs}
\usepackage[T1]{fontenc}
\usepackage{lmodern,microtype}
\usepackage[colorlinks=true,linkcolor=blue,urlcolor=blue]{hyperref}
\newtheorem{theorem}{Theorem}[section]
\newtheorem{lemma}[theorem]{Lemma}
\newtheorem{proposition}[theorem]{Proposition}
\newtheorem{corollary}[theorem]{Corollary}
\theoremstyle{remark}
\newtheorem{remark}[theorem]{Remark}
\newcommand{\norm}[1]{\left\|#1\right\|}
\title{Toward a four-dimensional Yang--Mills construction\\
\large Notebook 28: a fixed analytic ball and high-degree stability}
\author{Research working document}
\date{17 October 2026}
\setlength{\emergencystretch}{5em}
\begin{document}
\maketitle
\begin{abstract}
The generic Taylor-at-zero transport estimate shrinks the background
domain at each use. For the composite macrostep of Notebook 27,
we avoid this loss by composing inherited activities around a
bounded local core and expanding only the distant tail. The
pre-gauge representative is small by an inverse macro ratio.
A local analytic map supplies strict joint background/fluctuation
families for the entire inherited insertion. Cauchy estimates in
the fluctuation scaling parameter preserve its small factor in
the anchored activity norm. Classical precision and base
interactions are substituted into the stationary background
before Gaussian integration, using their independent classical
domains and finite support reserves. The Gaussian output is then
already a function of the next background.
A fixed sequence of parameter choices yields one inherited
small-field update on the same positive input and output
analytic radius, uniformly in preceding refinement and volume.
The general activity bound is not contractive. A separate
estimate proves an inverse-ratio contraction for background
transport when every incoming coefficient has vanishing
field Taylor jets through degree four. A bounded Taylor
subtraction produces such a remainder. The same estimate,
combined with a two-input Gaussian logarithm comparison,
contracts the full quantum response to input differences
whose coefficientwise low jets agree. Its high-degree
output projection can have Lipschitz constant at most
one-half. The full update can still regenerate low jets,
and this subtraction does not supply a gauge-invariant
extraction of them. These results close the stated domain
and support interfaces of a finite step and supply a
conditional quantum contraction, while leaving uniform
quantum amplitudes, coupling flow and the remaining
sectors of the existence argument open.
\end{abstract}

All identities below are first finite-regulator identities on
common small domains. Enlarged-domain identities are identified
as analytic continuations of those germs. This does not by itself
identify a globally continued homogeneous sector with the full
Wilson measure or control histories containing large-field
regions. The primitive ratio, composite blocking map, gauge
choice, retained link/color cutoff and normalization are those
constructed in Notebook 27.


% No preamble. Common-domain joint transport, unique prefix cdj-.
\section{Inherited interactions on a common background ball}
\label{cdj-section}

This section proves an anchored estimate for the actual inherited
insertion of the composite macrostep. It expands only a supported
tail around a bounded local core. The output field domain can be
the same product ball as the input domain. This does not assert
contraction of activity amplitudes or a complete quantum induction.
All scalar identities are on finite compatible periodic regulators;
all displayed anchored bounds are uniform in volume and in the
preceding classical refinement $M$.

\subsection{An abstract joint-support composition estimate}
\label{cdj-abstract-section}

Let $\Lambda_Q$ be the old carrier graph and $\Lambda_G$ the
common Gaussian and output-background graph. A projection
$\pi:\Lambda_Q\to\Lambda_G$ has fibers of size at most $F$,
takes connected carriers to connected sets, and has the stated
compression bound when used below. Let
\begin{equation}\label{cdj-input}
 \sup_i\sum_{H\ni i}e^{a_0|H|}\|E_H\|_b\le\epsilon,
\end{equation}
where $H$ is connected and $E_H$ is holomorphic in at most
$D|H|$ scalar current-field coordinates on the radius-$b$
product polydisc. Write $e_H=\|E_H\|_b$.

The output $w$ lies in a fixed complex product ball, possibly
of the same radius $b$. The retained fluctuation $\xi$ lies in
its actual product of real link/color balls. For
$|\lambda|<L_*$ suppose the field substituted in $E_H$ has
a decomposition
\begin{equation}\label{cdj-core-tail}
 f_\lambda|_H=c_{H,\lambda}+v_{H,\lambda}.
\end{equation}
The core is holomorphic in $(w,\lambda)$, measurable in $\xi$,
and reads both kinds of variables only in a connected carrier
$K_0(H)$ containing $\pi H$, where
\begin{equation}\label{cdj-core-hypotheses}
 |K_0(H)|\le c_0+c_1|H|,\qquad
 K_0(H)\subset N_{h_0}(\pi H),\qquad
 \|c_{H,\lambda}\|_\infty\le b-m,\quad 0<m<b.
\end{equation}
For each old scalar coordinate $(i,\alpha)$ let
\begin{equation}\label{cdj-tail-family}
 v_{H,\lambda;i,\alpha}
   =\sum_Xv_{H,\lambda;i,\alpha;X}
                     (w_X,\xi|_{K_0(H)}),\qquad
 \sup_{H,i,\alpha}\sum_X
 e^{\kappa\tau(\{\pi i\}\cup X)}
       \|v_{H,\lambda;i,\alpha;X}\|\le s.
\end{equation}
The coefficients are strict in their $w$ support $X$.
Their $\xi$ support stays inside the same $K_0(H)$,
even when $X$ is remote. The empty $X$ is allowed, and
no vanishing at $w=0$ is required. A stronger complete-fiber
norm, with the coordinate maximum inside the support sum,
is a sufficient input. All support assignments are fixed
independently of $\lambda$. Supremum norms include the
retained window and smaller closed $\lambda$ discs.
Here $\tau$ is graph Steiner length.

\begin{theorem}[Joint core/tail transport]\label{cdj-transport}
Put $t=s/m$ and suppose, for $B\ge0$ and $\delta>0$,
\begin{equation}\label{cdj-budget}
 t<1,\qquad B+\delta\le\kappa,\qquad
 c_1B+D\log\frac1{1-t}\le a_0.
\end{equation}
There are connected strict joint carriers $K$ with coefficients
$U_{K,\lambda}(w_K,\xi_K)$ such that
\begin{align}
 \sum_HE_H(f_\lambda|_H)&=\sum_KU_{K,\lambda},\notag\\
 \|U_\lambda\|_B
 :=\sup_z\sum_{K\ni z}e^{B|K|}\|U_{K,\lambda}\|
 &\le C_T\epsilon,\qquad
 C_T=e^{Bc_0+\delta h_0}F S_\delta,                     \label{cdj-transport-bound}\\
 S_\delta&=\sup_z\sum_y e^{-\delta d(z,y)}<\infty.\notag
\end{align}
The identities and holomorphy follow by normally convergent
series on smaller closed parameter domains.
\end{theorem}

\begin{proof}
Expand only in the tail, at the varying local core:
\[
 E_H(c_{H,\lambda}+v_{H,\lambda})
  =\sum_{\nu\in\mathbb N^{D|H|}}
     c_{H,\nu,\lambda}(w,\xi)v_{H,\lambda}^{\,\nu},
 \qquad
 \|c_{H,\nu,\lambda}\|\le e_Hm^{-|\nu|}.
\]
This is scalar-coordinate product Cauchy about a center with
the margin in \eqref{cdj-core-hypotheses}. The coefficient reads
only $K_0(H)$. No Taylor expansion of the core is made.
If needed take Cauchy radii below $m$ and pass to the limit.

Substitute \eqref{cdj-tail-family} in each scalar occurrence,
and join its support $X$ to its root $\pi i$ by a deterministic
minimizing Steiner graph. Let $T$ be the sum of their lengths
and $K$ their union with $K_0(H)$. Every $\xi$ dependency
already lies in that core carrier. Hence $K$ is a strict
connected joint carrier and
\[
 |K|\le c_0+c_1|H|+T,\qquad
 z\in K\ \Longrightarrow\ d(z,\pi H)\le h_0+T.
\]
The anchored weight is bounded by
\[
 \mathbf1_{\{z\in K\}}e^{B|K|}
 \le e^{Bc_0+\delta h_0}e^{-\delta d(z,\pi H)}
      e^{Bc_1|H|}
      \prod_{\rm occurrences}e^{(B+\delta)\tau}.
\]
The absolute sum of support choices and scalar multiindices
at fixed $H$ therefore costs exactly
\[
 e_H e^{Bc_1|H|}\sum_\nu t^{|\nu|}
   =e_He^{[Bc_1+D\log(1/(1-t))]|H|}.
\]
Repeated support choices already contain their multiplicities;
no additional factorial is inserted. Finally
$e^{-\delta d(z,\pi H)}
\le\sum_{i\in H}e^{-\delta d(z,\pi i)}$.
Use \eqref{cdj-input}, the fiber bound and
\eqref{cdj-budget}. This gives \eqref{cdj-transport-bound},
including anchors in the core halo or on added trees.
The same positive majorant proves normal convergence.
\end{proof}

\subsection{The fluctuation increment retains its small factor}
\label{cdj-increment-section}

Suppose $L_*\ge2$ and $f_0$ is independent of $\xi$. Use the
same canonical support assignment throughout the $\lambda$ disc
and set $Q_K=U_{K,1}-U_{K,0}$. Then
\begin{equation}\label{cdj-Cauchy-increment}
 \sum_KQ_K=\sum_H[E_H(f_1|_H)-E_H(f_0|_H)],
 \qquad
 \|Q\|_B\le \frac{C_T\epsilon}{L_*-1}.
\end{equation}
For $1<\rho<L_*$ take the Cauchy integral in the anchored
activity Banach space. The $n$th coefficient norm is at most
$C_T\epsilon\rho^{-n}$. Equivalently, keep the Cauchy integral
before the positive anchored sum and bound the whole anchored
norm of its integrand. One must not replace this argument
by a sum of separate carrier suprema over $\lambda$.
Let $\rho\uparrow L_*$ and sum $n\ge1$ at $\lambda=1$.

The total at $\lambda=0$ is the inherited background
functional, not a vacuum merely because it is independent
of fluctuations. In the actual construction below its
coefficients can themselves be chosen $\xi$ independent.
If
$q=C_T\epsilon/(L_*-1)$, each $\|Q_K\|\le q$ and
\begin{equation}\label{cdj-multiplier}
 \sup_z\sum_{K\ni z}e^{B|K|}
                     \|e^{Q_K}-1\|\le e^q q.
\end{equation}
The Gaussian multiplier is the bounded extension
$\chi_K(\xi)(e^{Q_K}-1)$, set to zero outside the original
retained window on $K$. On the global window this is
exact. Additional indicators are retained separately.
For the strict Gaussian theorem it is enough to choose
\begin{equation}\label{cdj-Gaussian-weight}
 B\ge4+\log\mathcal K_\theta+\kappa_G,
\end{equation}
since a connected $K$ has metric spanning-tree length
$|K|-1$. Under the full precision, baseline, tail and inherited
budgets, complex interpolation of the inherited exponent gives
a connected-log difference bounded by a fixed constant times $q$.
The integrated theorem below supplies this argument with the
base exponent retained. This application also requires the
precision and all other multipliers to have their strict $w$
supports on this graph.

\subsection{The actual local composite fluctuation map}
\label{cdj-local-map-section}

Fix the primitive $\ell$. On the final unit graph, an output
fiber consists of all current link potentials with initial
sites in one $R$-block, with their maximum norm. Its dimension
is of order $dR^d\dim\mathfrak g$; it is not given an
$\ell^1$ coordinate norm. For every prescribed finite reserve
$A$, the supported pre-gauge theorem gives
\begin{equation}\label{cdj-ap-family}
 a_{p,i}(w)=\sum_Xa_{i;X}(w_X),\qquad
 \sup_i\sum_Xe^{\kappa_0\tau(\{i\}\cup X)+A|X|}
                       \|a_{i;X}\|_b\le C_A b/R.
\end{equation}
The complete-fiber maximum is inside the support sum.
$\kappa_0,C_A$ and its domain are fixed at the primitive
ratio before choosing $R=\ell^s$, uniformly in $s,k$.
Every coefficient vanishes at $w=0$.

For an independent small current potential $a$, define
\[
 V(a)=e^{ia},\qquad
 u(a;y)=1,\quad u(a;x)=P_R(V(a);y,x),\qquad
 V^\star(a)=V(a)^{u(a)}.
\]
Using the actual composite chart, put
\begin{align}
 \Omega_R(a,\zeta)(b)
  &=\frac1i\log\!\left[
    e^{\,i\operatorname{Ad}_{u(a;b_-)^{-1}}
                 \Phi_{V^\star(a)}(\zeta)(b)}
       e^{ia(b)}\right],\notag\\
 \Delta_R(a,\zeta)&=\Omega_R(a,\zeta)-a,\qquad
 \Delta_R(a,0)=0.                                      \label{cdj-actual-map}
\end{align}
This is defined for independent $a$ on a fixed
$R$-dependent product ball. Stationarity is only used when
$a=a_p(w)$ to identify the physical branch. The map itself
is an ordinary local analytic map.

A nonpivot output reads its own retained variable and
the local endpoint gauge. A pivot output reads retained
variables and backgrounds in its two endpoint $R$-blocks;
all gauges on those links are internal to the same blocks.
Thus the whole output fiber at $i$ reads $a,\zeta$ in a
fixed neighborhood $\mathcal N(i)$. The conservative choice
$\mathcal N(i)=B_2(i)$ contains every needed block.
Let $N_0=\sup_i|\mathcal N(i)|$ and $r_0=2$, counting
periodic multiplicities if necessary.

Finite-dimensional holomorphy, on a slightly larger local
domain, and the zero in \eqref{cdj-actual-map} imply
\begin{equation}\label{cdj-local-bound}
 \|\Delta_R(a,\zeta)_i\|
    \le C_R^{\rm loc}\|\zeta|_{\mathcal N(i)}\|_\infty
 \quad\text{for}\quad
 \|a|_{\mathcal N(i)}\|_\infty,\,
 \|\zeta|_{\mathcal N(i)}\|_\infty<\rho_R.
\end{equation}
Integrate the bounded $\zeta$ derivative along $t\zeta$
to obtain the linear factor. The complete output maximum
is retained. Constants may depend on $R$, but the map
contains no variables at the preceding mesh $1/M$.

\subsection{A supported family without summing fine coordinates}
\label{cdj-Banach-Taylor-section}

Expand only the local Banach background input:
\[
 \Delta_R(a,\zeta)_i
   =\sum_{n\ge0}V_{i,n}(\zeta)[a,\ldots,a],
 \qquad V_{i,n}=D_a^n\Delta_R(0,\zeta)_i/n!.
\]
Radial Cauchy bounds the diagonal homogeneous polynomial;
complex polarization costs $n^n/n!\le e^n$. After a
harmless radius shrink,
\begin{equation}\label{cdj-vertex-bound}
 \|V_{i,n}(\zeta)\|_{\rm multilinear}
       \le C_R^{\rm loc}\|\zeta\|(e/\rho_R)^n.
\end{equation}
This estimates whole Banach fibers and does not sum their
$R^d$ scalar directions. Each coefficient is holomorphic
in $\zeta$ and zero at $\zeta=0$.

Decompose each input slot over its at most $N_0$ neighboring
fibers and insert \eqref{cdj-ap-family}. Assign the union
of the child supports as the new $w$ support. Intermediate
roots are Steiner vertices, so
\begin{align}
 \tau(\{i\}\cup\textstyle\bigcup_lX_l)
 &\le\sum_l[\tau(\{j_l\}\cup X_l)+d(i,j_l)],\notag\\
 |\textstyle\bigcup_lX_l|&\le\sum_l|X_l|.                \label{cdj-union}
\end{align}
The degree-$n$ absolute rooted norm is therefore bounded by
\begin{equation}\label{cdj-degree-majorant}
 C_R^{\rm loc}\|\zeta\|
 \left[
  \frac{eN_0e^{\kappa_0r_0}C_A b}{R\rho_R}
 \right]^n.
\end{equation}
Choose $b$ so that the bracket is at most $1/2$.
The geometric sum proves a strict family
\begin{align}
 \Delta_R(a_p(w),\zeta)_i
   &=\sum_Xd_{i;X}(w_X,\zeta|_{\mathcal N(i)}),\notag\\
 \sup_i\sum_Xe^{\kappa_0\tau(\{i\}\cup X)+A|X|}
                        \|d_{i;X}\|
   &\le C_R\|\zeta\|.                                  \label{cdj-joint-family}
\end{align}
This keeps the coordinate maximum inside the support sum
and bounds each degree before summing; no supremum/order
interchange is used. The empty support coefficient
$V_{i,0}(\zeta)$ is retained. It is a local fluctuation
function independent of $w$, not a vacuum constant.
Every $d_{i;X}$ vanishes at $\zeta=0$.

For the actual retained threshold $T$ set
\begin{equation}\label{cdj-lambda-field}
 f_\lambda(w,\xi)=\Omega_R(a_p(w),\lambda g\xi).
\end{equation}
Its background and fluctuation families obey
\begin{equation}\label{cdj-full-family}
 \|a_p\|_{\kappa_0,A;b}\le C_A b/R,\qquad
 \|f_\lambda-a_p\|_{\kappa_0,A;b,T}
                      \le C_R|\lambda|\,|g|T.
\end{equation}
Every output-root coefficient reads $\xi$ only in
$\mathcal N(i)$, regardless of how far its $w$ support
extends. Choose
\begin{equation}\label{cdj-lambda-radius}
 R\ge8C_A,\qquad
 L_*=\frac{b}{8C_R|g|T}\ge2,\qquad L_*|g|T<\rho_R.
\end{equation}
The last condition is another fixed upper bound on $b$.
On this disc the absolute family norm of $f_\lambda$ is
at most $b/4$. Thus any truncation by its support labels
has component amplitude at most $b/4$. At $\lambda=0$
the family is exactly that of $a_p$, independent of $\xi$.

\subsection{The core halo and the tail after auxiliary grouping}
\label{cdj-grouping-section}

Set $\kappa_t=\kappa_0/2$. Put the labels
$X\subset B_{\ell_{\rm tail}}(i)$ in the core, including
the empty support, and all other labels in the tail.
Every tail label has
$\tau(\{i\}\cup X)>\ell_{\rm tail}$, hence
\begin{equation}\label{cdj-unit-tail}
 \sup_i\sum_{\rm tail\ X}
 e^{\kappa_t\tau(\{i\}\cup X)+A|X|}
                  \|f_{\lambda;i;X}\|
 \le(b/4)e^{-(\kappa_0/2)\ell_{\rm tail}}.
\end{equation}
Both families still read $\xi$ only in $\mathcal N(i)$.
The core's actual component amplitude is at most $b/4$.

Group the source/final-unit cells into centered odd $J$-blocks,
conveniently $J=\ell^j$. These Gaussian macrocells have
side $JR$ in current units and become $J$-blocks of the next
current lattice. Old quantum carriers are $J$-blocks of
the current lattice. Their projection $\pi_R$ to Gaussian
macrocells has ratio $R$ and fiber $R^d$.
Both tilings refine the common centered $JR$ tiling.

For any desired fixed macro tail exponent $\kappa$,
\[
 \tau_{\rm macro}(\{\rho_J i\}\cup\rho_JX)
 \le2J^{-1}\tau_{\rm unit}(\{i\}\cup X)+d|X|.
\]
Thus $2\kappa/J\le\kappa_0/2$ and $A\ge d\kappa$
pay the projection. Taking the complete macro output
maximum inside the support sum costs at most $J^d$
old output roots, giving
\begin{equation}\label{cdj-macro-tail}
 s_{\rm tail}\le J^d(b/4)e^{-(\kappa_0/2)\ell_{\rm tail}}.
\end{equation}
Each old $J$-current input fiber is a restriction of
the containing whole $JR$ macro output fiber, so it
inherits this bound at its projected macro root.
There is no further sum over its scalar coordinates here.

If $\ell_{\rm tail}+r_0\le J$, a conservative joint core
carrier for a connected old $H$ is
\begin{equation}\label{cdj-actual-core}
 K_0(H)=N_{\infty,2}(\pi_RH).
\end{equation}
It is connected and contains every core $w$ dependency
and every $\xi$ dependency of every tail term for outputs
in $H$. In particular
\[
 |K_0(H)|\le5^d|\pi_RH|,\qquad
 K_0(H)\subset N_{2d}(\pi_RH).
\]
Declaring the whole halo only enlarges the support.
The distant background labels do not bring distant
fluctuation dependencies.

Connected-carrier compression at odd $R\ge5$ gives
\[
 |\pi_RH|\le2^d+(2^d/k_R)|H|,
 \qquad k_R=1+\lfloor(R-3)/4\rfloor.
\]
The abstract theorem therefore uses
\begin{equation}\label{cdj-actual-constants}
 c_0=5^d2^d,\quad c_1=\frac{5^d2^d}{k_R},
 \quad h_0=2d,\quad F=R^d,\quad m=b/2.
\end{equation}
Its scalar-coordinate count is the actual old quantum-cell
value
\begin{equation}\label{cdj-D}
 D(J)=dJ^d\dim\mathfrak g,
\end{equation}
not the retained Gaussian dimension.

\subsection{A noncircular choice of parameters on the same ball}
\label{cdj-parameters-section}

Fix $B$ using a dimension-independent upper bound in
\eqref{cdj-Gaussian-weight}, before choosing $R,J$.
For example $\kappa_G=2$, $m_G=3$ and the eventual choice
$\theta=1/(4D_G)$ give $\log\mathcal K_\theta\le5/8$
for every $D_G$, so $B=53/8$ can be fixed now.
The later Gaussian dimension changes $\theta$, not $B$.
Fix $\delta>0$, choose $\kappa\ge B+\delta$, and choose
the finite reserve $A\ge d\kappa$ in
\eqref{cdj-ap-family}. Its constants are then fixed
at primitive $\ell$.

Choose $R=\ell^s$ large enough that
\begin{equation}\label{cdj-R-choice}
 R\ge8C_A,\qquad
 \frac{5^d2^d}{k_R}B\le a_0/2.
\end{equation}
Choose $J=\ell^j$ large enough for the above spatial
projection and the separate classical/Gaussian regrouping
conditions. With
$\ell_{\rm tail}=c\log J+O(1)$, where $c>4d/\kappa_0$,
one has $\ell_{\rm tail}+r_0\le J$ eventually, while
\begin{equation}\label{cdj-dimension-tail}
 D(J)\frac{s_{\rm tail}}m
 \le\frac{d\dim\mathfrak g}{2}J^{2d}
          e^{-(\kappa_0/2)\ell_{\rm tail}}\longrightarrow0.
\end{equation}
Increase $J$ and the fixed offset so that
$s_{\rm tail}/m\le1/2$ and
$D(J)\log[1/(1-s_{\rm tail}/m)]\le a_0/2$.
These choices are independent of $b$: it cancels from
the tail-to-margin ratio.

Now choose $b$ once below all classical and chart radii,
including \eqref{cdj-degree-majorant} and
\eqref{cdj-lambda-radius}. Finally choose $|g|T$ for
$L_*\ge2$ and the desired interaction budgets.
The same radius $b$ is the incoming inherited domain
and the outgoing background domain. It has not been
replaced by a smaller recursive $b/C$ at this step.

The actual map now supplies all the hypotheses of
Theorem~\ref{cdj-transport}. Equations
\eqref{cdj-Cauchy-increment} and
\eqref{cdj-lambda-radius} yield the strict joint bound
\begin{equation}\label{cdj-actual-increment}
 \|U_1-U_0\|_B
 \le\frac{C_T\epsilon}{L_*-1}
 \le16 C_T C_R\frac{|g|T}{b}\epsilon,\qquad
 C_T=e^{B5^d2^d+2d\delta}R^dS_\delta.
\end{equation}
The exact inherited multiplier norm is at most $e^q q$
with $q$ the first bound in
\eqref{cdj-actual-increment}. The total $U_0$ is
$E(a_p(w))$, with the same stronger carrier bound
$C_T\epsilon$, and can be chosen termwise independent
of $\xi$. It has no $g$ factor and is retained as a
background activity, not discarded as a vacuum.

\subsection{Scope of the common-domain estimate}
\label{cdj-scope}

The finite-regulator identity and gauge comparison are
those of the composite integral module: on the common
overlap total gauge invariance identifies
$f_\lambda$ with the actual rooted insertion, while
the pre-gauge expression defines its continuation on
the common $b$-ball. No coefficientwise gauge invariance,
fine-coordinate $\ell^1$ count, or rooted $1/R$ angle
contraction is assumed.

The precision and purely classical base multipliers
must also be supplied with strict background/joint
supports on this graph before applying the Gaussian
theorem directly in $w$. They have independent classical
domains and arbitrary finite classical support reserves;
their interface is a separate module. Uniform complex
nonvanishing follows only after all cluster hypotheses
and budgets have been met.

The constants here may be extremely large.
The background term has a finite prefactor, not a
contraction, and only the fluctuation difference gains
$|g|T/b$. Vacuum and marginal extraction, coupling flow,
large-field complements, mixed histories and observable
bounds remain additional requirements for an existence
proof.

\section{Classical support substitution before Gaussian integration}
\label{cdc-section}

This section places the classical precision and base factors directly
in the next background $w$, before applying the Gaussian expansion.
The composite ratio $R=\ell^s$ is fixed. Constants may depend on
$R$ and on the later auxiliary grouping, but are uniform in the
preceding homogeneous depth $M$ and compatible volume. The inputs
are the established classical strict families with arbitrary fixed
cardinality reserve and the actual supported stationary map. No
such arbitrary reserve is assumed for inherited quantum activities.

\subsection{A precision composition lemma}
\label{cdc-precision-section}

On one fixed cell graph let
$B(a)=\sum_Y B_Y(a|_Y)$ be a symmetric matrix family. Its matrix
rows and columns are supported in $Y$, and $B_Y(0)=0$. Each
$a$-cell has $D$ scalar coordinates with independent analytic
radius $b_{\rm cl}$. Write $\tau(Y)$ for graph Steiner-tree length,
allowing auxiliary vertices. Suppose
\begin{equation}
 \sup_i\sum_{Y\ni i}\sum_{j\in Y}
 e^{\sigma\tau(Y)+A|Y|}
       \|(B_Y)_{ij}\|_{b_{\rm cl}}\leq N_B.
 \label{cdc-precision-input}
\end{equation}
The norms inside this sum are block operator norms. Let the
stationary substitution have a normally convergent strict family
\begin{equation}
 a_i(w)=\sum_X u_{i;X}(w_X),\qquad u_{i;X}(0)=0,\qquad
 \sup_i\sum_Xe^{\kappa\tau(\{i\}\cup X)+a|X|}
                         \|u_{i;X}\|_r\leq s(r).
 \label{cdc-map-input}
\end{equation}
Here each output norm is the maximum over the entire cell fiber,
\emph{inside} the support sum.

\begin{lemma}[Composition in the matrix row norm]
\label{cdc-precision-lemma}
Choose $0<\kappa'<\min(\sigma,\kappa)$, $A'\geq0$ with
$a\geq A'$, and $\rho>0$. Put
\begin{equation}
 t=\frac{s(r)}{b_{\rm cl}}<1,\qquad
 \alpha_D(t)=D\log\frac1{1-t}.
 \label{cdc-alpha}
\end{equation}
If $A'+\alpha_D(t)+\rho\leq A$, then
$B(a(w))=\sum_Z B'_Z(w_Z)$ has symmetric strict pieces with
rows and columns in $Z$, and
\begin{equation}
 \sup_i\sum_{Z\ni i}\sum_{j\in Z}
 e^{\kappa'\tau(Z)+A'|Z|}\|(B'_Z)_{ij}\|_r
 \leq N_B\alpha_D(t)
                   \sup_{n\geq1}n e^{-\rho n}.
 \label{cdc-precision-output}
\end{equation}
In particular the bound is $O(r)$ if $s(r)=O(r)$ and the
displayed margins are fixed.
\end{lemma}
\begin{proof}
Expand a matrix entry of $B_Y$ on its scalar-coordinate polydisc.
Cauchy bounds its multiindex coefficient by
$\|(B_Y)_{ij}\|_{b_{\rm cl}}b_{\rm cl}^{-|\nu|}$.
The zero multiindex is absent. Substitute the strict family
\eqref{cdc-map-input} in every occurrence of every scalar input.
If an occurrence is rooted at $k\in Y$ and assigned support $X$,
declare the resulting carrier to be
$Z=Y\cup\bigcup X$. Steiner gluing gives
\begin{equation}
 \tau(Z)\leq\tau(Y)+\sum_{\rm occurrences}
                            \tau(\{k\}\cup X),\qquad
 |Z|\leq|Y|+\sum_{\rm occurrences}|X|.
 \label{cdc-precision-gluing}
\end{equation}
The scalar multiindex sum is
\begin{equation}
 (1-t)^{-D|Y|}-1
 \leq\alpha_D(t)|Y|e^{\alpha_D(t)|Y|}.
 \label{cdc-precision-generating}
\end{equation}
The $\rho$ reserve pays $|Y|$ and proves
\eqref{cdc-precision-output}. For a fixed nonzero matrix row
$i$, the original piece already had $i\in Y$, and its nonzero
columns remain in $Y$. Substitution does not create new matrix
rows or columns. Thus no path-anchor rerooting factor is needed.
Symmetry is preserved coefficientwise and after grouping equal
$Z$. Absolute convergence justifies all groupings.
\end{proof}

If only $\sum_YB_Y(0)=0$ is initially known, subtract the
individual constants. Their absolutely convergent sum is zero,
and the input norm increases by at most a factor two. The lemma
is a statement about the precision's weighted row norm. It does
not assert a norm of whole global matrices anchored at every
new support site.

\subsection{Application to the actual composite precision}
\label{cdc-actual-precision-section}

Group current-unit classical families into source $R$-blocks.
The number of current background coordinates in each fiber is
finite at fixed $R$. The established classical support
construction, followed by the exact finite-range composite pivot
and ordinary logarithm maps, supplies individually symmetric
strict pieces of
\begin{equation}
 B(a)=A_{M,R}(a)-A_{M,R,0}
 \label{cdc-actual-precision}
\end{equation}
on an independent current-background product ball. Their row
norm has an arbitrary prescribed finite cardinality reserve and
some positive spatial exponent, uniformly in $M$ and volume.
The chain rule includes the classical current contracted against
both chart-curvature terms, and the full absolute gauge Hessian.
These terms are not omitted away from the physical branch.

The actual stationary field is the pre-gauge supported
representative followed by the fixed-$R$ internal contour gauge.
Finite-range Banach-valued analytic composition therefore gives
\eqref{cdc-map-input} for $a_*(w)$ with any fixed finite
cardinality reserve, a positive exponent and $s(r)=O_R(r)$.
Its output fiber is the entire current field in an $R$-block,
with its maximum norm inside the support sum. The lemma gives
a symmetric strict $w$-precision with row norm $O_R(r)$.
The outer radius $b_{\rm cl}$ here is a classical radius; it is
not the inherited quantum radius reused at every step.

Now group source roots and next-background cells into odd
$J$-blocks. For a projected support $\varrho_J Z$, the established
geometric bounds give
\begin{equation}
 \tau_{\rm mac}(\varrho_JZ)
       \leq\frac2J\tau(Z)+d|Z|,\qquad
 \ell_{\rm mac}(\varrho_JZ)
       \leq\frac4J\tau(Z)+2d|Z|.
 \label{cdc-projection}
\end{equation}
The second uses metric spanning-tree length $\ell$, and is
conservative. Grouping matrix rows costs at most $J^d$:
sum the old rows in the macrocell before taking its block norm.
For the target strict Gaussian exponent $\gamma=8$, it is
sufficient in $d=4$ to choose
\begin{equation}
 \frac{32}{J}\leq\kappa',\qquad A'\geq64.
 \label{cdc-precision-reserve}
\end{equation}
This produces the required macro precision norm $O_{R,J}(r)$.
The separately established covariance decay is regrouped at the
same $J$; increasing $J$ if necessary supplies the fixed
Gaussian spatial exponents.

For clarity, the small-radius choice is made \emph{after} $R,J$.
The real baseline retains its spectral bounds
$\gamma_R I\leq A_{M,R,0}\leq B_RI$ under grouping. If
$\nu(r)$ bounds the substituted precision perturbation in the
macro weighted row norm, then
\begin{equation}
 \nu(r)\leq\gamma_R/2
 \quad\Longrightarrow\quad
 \|A(w)^{-1}-A_0^{-1}\|
             \leq2\gamma_R^{-2}\nu(r).
 \label{cdc-covariance-comparison}
\end{equation}
Hence the sufficient additional condition
\begin{equation}
 \nu(r)\leq\frac{\theta\gamma_R^2}{4B_R}
 \label{cdc-relative-neighborhood}
\end{equation}
puts the covariance within relative norm $\theta/2$ of the real
baseline, leaving the other half for the Gaussian Cauchy
perturbation. The strict density-support construction has
additional positive smallness thresholds in this same weighted
precision norm. Since $\nu(r)=O_{R,J}(r)$, all can be met by
choosing a single final $b\leq r$ after the finite geometry and
$\theta$ have been fixed. For example
$\theta=1/(4D_G)$ uses the actual finite macro Gaussian
dimension $D_G$. This choice does not require replacing the
inherited field radius by a smaller fraction after each
stationary substitution.

\subsection{Scalar exponents with spectator fluctuation variables}
\label{cdc-scalar-section}

Let $V_Y(a|_Y,\xi|_Y)$ be a classical or measure exponent family,
holomorphic in the independent background, measurable on the
actual retained cutoff, with
\begin{equation}
 \sup_i\sum_{Y\ni i}e^{\sigma\tau(Y)+A|Y|}
             \|V_Y\|_{b_{\rm cl},{\rm cutoff}}
 \leq\epsilon_{\rm cl}.
 \label{cdc-scalar-input}
\end{equation}
Only $Y$ is read in the fluctuation variables. It need not be
connected. Empty scalar vacuum terms, if present, are tracked
separately. Substitute \eqref{cdc-map-input} in the background
slots while keeping $\xi$ as spectator data.

\begin{lemma}[Joint macro carriers for the base factors]
\label{cdc-scalar-lemma}
Fix a target cardinality exponent $\mathsf B\geq0$ and
$\delta>0$. With $t=s(r)/b_{\rm cl}<1$, suppose
\begin{equation}
 \frac{2(\mathsf B+\delta)}J\leq\min(\sigma,\kappa),\qquad
 a\geq d(\mathsf B+\delta),\qquad
 A\geq d(\mathsf B+\delta)+\alpha_D(t).
 \label{cdc-scalar-budget}
\end{equation}
There is a normally convergent decomposition into connected
strict joint $(w,\xi)$ macro carriers $K$ such that
\begin{equation}
 \sum_YV_Y(a_*(w)|_Y,\xi|_Y)=\sum_KV'_K(w_K,\xi_K),
 \qquad
 \sup_z\sum_{K\ni z}e^{\mathsf B|K|}
                         \|V'_K\|
 \leq e^{\mathsf B}J^d S_\delta\epsilon_{\rm cl},
 \label{cdc-scalar-output}
\end{equation}
where $S_\delta=\sup_z\sum_y e^{-\delta d(z,y)}$.
\end{lemma}
\begin{proof}
Taylor expand in the $D|Y|$ scalar background coordinates only.
Retain $Y$ as a fluctuation carrier even in the zero Taylor
coefficient. Join $\varrho_JY$ by a deterministic Steiner graph;
for each substituted support $X$ rooted at $k\in Y$, attach a
deterministic Steiner graph spanning
$\{\varrho_Jk\}\cup\varrho_JX$. Their union is a connected
joint carrier $K$. If $T$ is the sum of their edge lengths, then
\begin{align}
 |K|&\leq1+T,\nonumber\\
 T&\leq\frac2J\tau(Y)+d|Y|
       +\sum_{\rm occurrences}
          \left[\frac2J\tau(\{k\}\cup X)+d|X|\right].
 \label{cdc-joint-tree}
\end{align}
The carrier includes every background argument and the entire
projected fluctuation carrier, regardless of their overlap.
For $z\in K$, $d(z,\varrho_JY)\leq T$. Thus
\begin{equation}
 e^{\mathsf B|K|}\mathbf1_{z\in K}
 \leq e^{\mathsf B}e^{-\delta d(z,\varrho_JY)}
                                      e^{(\mathsf B+\delta)T}.
 \label{cdc-reroot}
\end{equation}
The spatial and cardinal weights in
\eqref{cdc-scalar-budget} pay the last factor. The scalar
multiindex sum, now including degree zero, is exactly
$(1-t)^{-D|Y|}=e^{\alpha_D(t)|Y|}$, paid by the remaining
input cardinal reserve. Finally use
\begin{equation}
 e^{-\delta d(z,\varrho_JY)}
 \leq\sum_{y\in\varrho_JY}e^{-\delta d(z,y)}.
 \label{cdc-anchor-sum}
\end{equation}
For each projected root $y$ there are at most $J^d$ source
sites. The original anchored norm and the lattice sum
$S_\delta$ yield \eqref{cdc-scalar-output}. These estimates
are positive majorants uniform over the actual real cutoff;
normal convergence gives background holomorphy. All deterministic
support assignments are fixed before choosing the fields.
\end{proof}

The zero background Taylor term can depend on $\xi$ and is then
an interaction, not a vacuum constant. The proof does not infer
a strict family from a bound on the summed scalar functional.
It starts with the actual coefficient bound
\eqref{cdc-scalar-input}.

\subsection{Actual classical and measure inputs and scope}
\label{cdc-matching-section}

At fixed $R$, the already established classical support theorem
applies to the action after the actual composite chart and
ordinary-log conversion. Its Taylor remainder after subtracting
constant, linear and quadratic terms has arbitrarily weighted
strict coefficients. On the physical stationary branch the
subtracted linear term vanishes, so this extension agrees with
the exact centered integral. The full absolute gauge action is
treated the same way; no off-stationary residual is discarded.

The local ambient pivot logarithms, Haar logarithms and gauge
remainder also have actual independent-background coefficient
families. Composite one-pivot incidence keeps their ranges
finite at fixed $R$. The retained link cutoff is unchanged.
With its link threshold denoted by $T$, these families retain
the respective small factors
\begin{equation}
 O_{R,J}\bigl(|g|T^3+|g|T+g^2T^2\bigr)
 \label{cdc-base-smallness}
\end{equation}
after the scalar substitution, provided the fixed chart and
Taylor domains hold. Equal-carrier exponents are grouped before
exponentiation. If their anchored exponent bound is $v$, it
bounds each individual exponent and gives the multiplier norm
bound $e^v v$ for $e^{V'_K}-1$.

Background-only pivot-delta factors are treated by the same
scalar substitution. The Gaussian determinant is expanded in
closed chains using the already substituted strict precision;
its field-independent normalization remains separate. The
classical outer radii and reserves are fixed independently of
the old quantum activity radius.

Thus the precision and base multipliers already read $w$ on
the auxiliary Gaussian graph. Once a separate inherited joint
family meets its budget, the Gaussian expansion outputs
$w$-activities directly on next-current $J$-cells. There is no
second transport through $a_*(w)$ after integration. This changes
only the order of normally convergent support substitution;
the exact integration variables, covariance, cutoff and
normalization are unchanged.

This section is not the integrated quantum-step theorem. It
does not estimate inherited activity amplitudes, prove a
contractive recursion, identify a running coupling or control
large fields. All equality statements start on their common
local analytic domains and extend only on the established
holomorphic continuations.

\section{One composite small-field step on the same analytic ball}
\label{cdm-section}

We combine Sections~\ref{cdj-section} and \ref{cdc-section} with the
actual composite constraint, Gaussian and integral of Notebook~27.
All statements are first finite-regulator statements on compatible
periodic tori. Constants are uniform in the preceding homogeneous
depth $M=\ell^k$ and volume, after the primitive odd ratio $\ell$,
macro ratio $R=\ell^s$ and auxiliary grouping $J=\ell^j$ are fixed.
The theorem concerns a local holomorphic small-field map. The measure
identification required beyond its original germ is stated explicitly.

Set $d=4$ and use the fixed trace-orthonormal real color coordinates.
The input activity graph consists of $J$-blocks of the current
unit lattice. The output graph consists of $J$-blocks of the next
lattice, equivalently $JR$-blocks of the current lattice, and is
also the Gaussian integration graph. Thus input and output use
the same grouping $J$ in their respective current units. Projection
$\pi_R$ has fiber $R^d$, preserves connected images, and satisfies
\begin{equation}
 |\pi_RH|\le 2^d+\frac{2^d}{k_R}|H|,
 \qquad k_R=1+\left\lfloor\frac{R-3}{4}\right\rfloor.
 \label{cdm-compression}
\end{equation}
The distinct lifted block interfaces are retained on the chosen
cofinal regulator family.

\subsection{Statement and exact field-dependent factors}
\label{cdm-statement-section}

Write the inherited total functional as
\begin{equation}
 E(a)=\sum_H E_H(a|_H),\qquad
 \norm{E}_{a_0;b}:=
 \sup_i\sum_{H\ni i}e^{a_0|H|}\norm{E_H}_b
 \le\epsilon\le\bar\epsilon,\qquad a_0=\frac1{16},
 \label{cdm-input}
\end{equation}
where the carriers are connected and each coefficient is holomorphic
on the scalar-coordinate product ball of radius $b$. The fixed
$\bar\epsilon<\infty$ bounds the input class. Assume local gauge
invariance of the \emph{total} $E$ on its original common domain.
If $E$ is a logarithm, its compatible holomorphic branch is part
of the input; coefficientwise gauge invariance is not assumed.

\begin{theorem}[A common-domain composite step]
\label{cdm-theorem}
There are finite choices $R,J$, positive $b,g_0$, and retained
link/color-ball thresholds $T=T(g)\ge1$ such that, for
$0<g<g_0$, the chosen small-field macrostep has a holomorphic
continuation to the entire next-angle ball $\norm{w}_\infty<b$.
With $W=e^{iw}$ it has the factorization
\begin{equation}
 \rho_{\mathrm{next}}(W)=
 N_{M,R,g}\exp\left[-g^{-2}S_{MR}(W)+E_{\mathrm{next}}(w)\right].
 \label{cdm-factorization}
\end{equation}
Here $N_{M,R,g}$ is field independent, and $S_{MR}$ is the
homogeneous ordinary classical minimum at total depth $MR$,
after the $d=4$ physical rescaling. No coupling renormalization
has yet been performed.

The output has connected strict $w$-activities with the same
angle radius $b$ and cardinality exponent $a_0$. In the notation
constructed below,
\begin{align}
 \norm{E_{\mathrm{next}}}_{a_0;b}
 &\le C_T\epsilon+(C_\Delta+C_{\det})b
       +2S_1\left[e^3\tau(T)+e^{v+q}(v+q)\right],
 \label{cdm-output-bound}\\
 v&\le C_{\mathrm{base}}(gT^3+gT+g^2T^2),\qquad
 q\le16C_TC_R\frac{gT}{b}\epsilon .
 \label{cdm-small-factors}
\end{align}
All constants are uniform in $M$ and compatible volume. The
Gaussian cutoff integral is nonzero on this whole ball and has
the analytic logarithm used in \eqref{cdm-factorization}.
The total correction satisfies the local gauge Ward identity
throughout the ball, and is invariant along gauge paths that
remain in it.

These are statements about the continuation of the actual local
small-field germ. They identify the exact chosen macro integral
throughout the ball when a gauge-orbit current density with the
stated pre-gauge representation is available. They do not identify
all quantum histories or the full Wilson pushforward merely from
holomorphic uniqueness.
\end{theorem}

The following construction proves the theorem and specifies its
constants and domains. The estimate is not an amplitude
contraction: $C_T$ can be much larger than one.

Use the actual composite ordinary map
$F_R=\mathcal M_\ell^s$ and the independent direct-$R$ contour
gauge. The physical backgrounds of Notebook~27 satisfy
\begin{equation}
 \begin{aligned} V_p(W)&=\mathcal M_\ell^k U_{MR}(W),\qquad
 V^\star(W)=V_p(W)^{u(W)},\\
P_R(V^\star)&=1,
 \qquad S_M(V^\star)=S_M(V_p)=S_{MR}(W).
 \end{aligned}
\label{cdm-background}
\end{equation}
The last identity follows from the two actual minimum comparisons.
Current spacing is one and the next spacing is $R$; the
coupling-free action has no extra rescaling factor in dimension
four.

The retained chart preserves its nonpivot coordinates exactly.
The integration field is $e^{i\Phi_{V^\star}(g\xi)}V^\star$ and
the window is
$\prod_{b\ {\rm retained}}\mathbf1_{\{|\xi(b)|\le T\}}$.
The coupling stays inside $\Phi$. Set
\begin{align}
 A(w)&=
 D_z^2\left[
 S_M(e^{i\Phi_{V^\star}(z)}V^\star)
 +\mathcal G_R(e^{i\Phi_{V^\star}(z)}V^\star)\right]_{z=0},
 \notag\\
 E_{\mathrm{pre}}(w)&=E(a_p(w)),\qquad a_p=\frac1i\log V_p,
 \notag\\
 Z(w)&=\mathbb E_{A(w)^{-1}}\left[
 \chi_{\mathrm{window}}(\xi)
 e^{V_{\mathrm{base}}(w,\xi)+Q_{\mathrm{inh}}(w,\xi)}
 \right].
 \label{cdm-integral-parts}
\end{align}
The base exponent is the centered classical/gauge remainder,
the ambient pivot-straightening logarithms, and the Haar
logarithms. The inherited difference is exactly
\begin{align}
 Q_{\mathrm{inh}}(w,\xi)&=E(f_1(w,\xi))-E(f_0(w)),\notag\\
 f_\lambda(w,\xi)(b)&=
 \frac1i\log\left[
 e^{\,i\operatorname{Ad}_{u(w;b_-)^{-1}}
                \Phi_{V^\star(w)}(\lambda g\xi)(b)}
 V_p(w)(b)\right],\qquad f_0=a_p.
 \label{cdm-pre-gauge-insertion}
\end{align}
This group identity changes no integration variable $\xi$.

For clarity, let $n_e$ count positive current bonds, $n_c$ final
macro bonds, $n_s$ noncenter current sites, and
$n_z=\dim\mathfrak g\,(n_e-n_c)$. In trace-orthonormal exponential
coordinates write $d(e^{iX})=h_0j(X)dX$, $j(0)=1$, and put
\[
 c(g)=\int_G
 \exp[-(1-\operatorname{Re}\tau(h))/g^2]\,dh.
\]
With $\mathcal D_R(V)$ the actual pivot-delta determinant and
$A_0=A_{M,R,0}$ the flat retained precision, the field-independent
normalization before any additional vacuum subtraction is
\begin{equation}
 N_{M,R,g}=
 c(g)^{-n_s}h_0^{\,n_e-n_c}g^{n_z}(2\pi)^{n_z/2}
       \mathcal D_R(1)\det(A_0)^{-1/2},
 \label{cdm-normalization}
\end{equation}
times any separately tracked old field-independent normalization.
The background-dependent parts of the pivot and Gaussian
determinants remain in $E_{\mathrm{next}}$.

On the common overlap this is precisely the local Haar/delta
integral. If an actual gauge-invariant current density on gauge
orbits has representation $e^{-S_M/g^2+E}$ at the pre-gauge
fields in \eqref{cdm-pre-gauge-insertion}, the same expression
is the exact selected sector pointwise throughout the ball.
Otherwise it specifies only the verified analytic continuation
of that sector's local germ. The measurable completion and
one-shot sector are those of Notebook~27, not an asserted
intersection of $s$ sequential primitive sectors. Additional
original indicators would have to be retained and estimated.

\subsection{An explicit order of parameter choices}
\label{cdm-parameters-section}

First fix the exponents
\begin{equation}
 a_0=\frac1{16},\quad \kappa_G=2,\quad m_G=3,\quad
 \gamma=8,\quad B=\frac{53}{8},\quad \delta=1,\quad
 \kappa_{\mathrm{tail}}=8.
 \label{cdm-exponents}
\end{equation}
The tuned Gaussian density factor below has logarithm at most
$5/8$, so \eqref{cdj-Gaussian-weight} holds, and
$\kappa_{\mathrm{tail}}\ge B+\delta$.

Choose a fixed pre-gauge cardinal reserve
$A_{\mathrm{in}}\ge d\kappa_{\mathrm{tail}}$ before $R$.
By \eqref{cdj-ap-family}, there are $C_A,\kappa_0,r_A>0$ fixed
at the primitive ratio, independently of $R,M$ and volume, with
\begin{equation}
 \sup_i\sum_Xe^{\kappa_0\tau(\{i\}\cup X)+A_{\mathrm{in}}|X|}
                  \norm{a_{p,i;X}}_r
       \le C_A r/R,\qquad 0<r\le r_A.
 \label{cdm-pregauge-family}
\end{equation}
The norm takes the maximum over the whole current-link fiber
of one $R$-block \emph{inside} the support sum.
Choose a sufficiently large power $R=\ell^s$ so that
\begin{equation}
 R\ge8C_A,\qquad
 \frac{5^d2^d}{k_R}B\le a_0/2.
 \label{cdm-R-choice}
\end{equation}
The composite chart and its local radii are now fixed. The actual
flat retained Gaussian has uniform bounds
\begin{equation}
 \gamma_RI\le A_0\le B_RI,\qquad
 \lambda_0=B_R^{-1},\quad \Lambda_0=\gamma_R^{-1},
 \label{cdm-spectral-bounds}
\end{equation}
and an exponentially local inverse. These constants may depend
on $R$ but not on the old depth $M$.

Next choose any further \emph{classical-only} reserves needed
for the fixed-$R$ chart and the two lemmas in
Section~\ref{cdc-section}. They may depend on $R$ and use
independent classical outer radii. They do not change
$A_{\mathrm{in}}$ or $C_A$ in
\eqref{cdm-pregauge-family}--\eqref{cdm-R-choice}.
The rooted stationary field is a fixed-$R$ local analytic
function of $a_p$; its support norm is $O_R(r)$ with any such
fixed reserve. Lemma~\ref{cdc-precision-lemma} therefore
provides strict symmetric precision pieces with row norm
$O_R(r)$.

Choose $J=\ell^j$ sufficiently large for the grouped covariance
and precision exponents, in particular
\begin{equation}
 \frac{32}{J}\le\kappa',\qquad A'\ge64,\qquad
 \frac{16}{J}\le\kappa_0/2,
 \label{cdm-J-choice}
\end{equation}
with the available $\kappa',A'$ in
\eqref{cdc-precision-reserve}. Set
$\ell_{\mathrm{tail}}=c\log J$ plus a fixed offset, with
$c>4d/\kappa_0$, so that eventually
$\ell_{\mathrm{tail}}+2\le J$. The actual joint family of
Section~\ref{cdj-grouping-section} then has
\begin{align}
 s_{\mathrm{tail}}
 &\le J^d(b/4)e^{-\kappa_0\ell_{\mathrm{tail}}/2},
 &m_{\mathrm{core}}&=b/2,\notag\\
 c_0&=5^d2^d,&c_1&=5^d2^d/k_R,\notag\\
 h_0^{\mathrm{halo}}&=2d,&F&=R^d,\qquad
 D(J)=d\,\dim\mathfrak g\,J^d .
 \label{cdm-tail-data}
\end{align}
Here $h_0^{\mathrm{halo}}$ is a graph radius, distinct from the
Haar constant in \eqref{cdm-normalization}. Increase $J$ and
the tail offset until, with $t=s_{\mathrm{tail}}/m_{\mathrm{core}}$,
\begin{equation}
 t\le\tfrac12,\qquad D(J)\log\frac1{1-t}\le a_0/2.
 \label{cdm-tail-budget}
\end{equation}
This is possible because its sufficient bound is a constant
times $J^{2d}e^{-\kappa_0\ell_{\mathrm{tail}}/2}\to0$.
The ratio $t$ in this sufficient estimate is independent of $b$.
The entire local $\xi$ collar and the core lie in the fixed
two-neighbor halo used in \eqref{cdj-actual-core}.

Only now choose the final field radius. The macro retained
dimension and tunable density parameters are
\begin{align}
 D_G&=J^d d(R^d-1)\dim\mathfrak g,\qquad
 \theta=\frac1{4D_G},\notag\\
 \mathcal K_\theta
 &=\left[\frac{1+\theta}{(1-\theta)^3}\right]^{D_G/2},
 &\log\mathcal K_\theta&\le\frac58,\qquad
 \rho_\theta=\frac{1+\theta}{1-\theta}.
 \label{cdm-density-parameters}
\end{align}
The uniform Gaussian bounds fix a positive common strictly
local threshold $\eta_*$. Shrink it if necessary so that
$\eta_*\le1$.

Choose $b>0$ below all classical, physical-branch and chart
radii. In particular require
\begin{equation}
 \frac{eN_0e^{2\kappa_0}C_A b}{R\rho_R}\le\frac12,
 \label{cdm-joint-domain}
\end{equation}
as in \eqref{cdj-degree-majorant}. The substituted macro
precision norm $\nu(b)=O_{R,J}(b)$ must meet every strictly
local support-algebra threshold. Include
\begin{equation}
 \nu(b)\le\gamma_R/2,\qquad
 \nu(b)\le\frac{\theta\gamma_R^2}{4B_R},
 \label{cdm-precision-small}
\end{equation}
as justified in
\eqref{cdc-covariance-comparison}--\eqref{cdc-relative-neighborhood}.
Also impose the tunable density-algebra conditions
$u_G\le\theta$ and
$w_G\le\theta/[(1+\theta)\Lambda_0]$, where $u_G,w_G$
are its continuous $O(\nu)$ support-algebra quantities, and
the smallness for the determinant chain expansion below.
All hold after one shrink of $b$ with $R,J$ fixed.
This selects a single input and output radius, rather than
dividing the inherited radius during substitution.

A bad macrocell for the actual link/color-ball window implies
that its full Euclidean retained norm exceeds $T$. Its valid
joint-tail constant is therefore
\begin{equation}
 \tau(T)=\mathcal K_\theta\,2^{D_G/2}
       \exp\left[-\frac{T^2}{4\rho_\theta\Lambda_0}\right].
 \label{cdm-tail-constant}
\end{equation}
Choose $T\ge1$ with $e^3\tau(T)\le\eta_*/4$.
This estimates the bad event without replacing the actual
good window by a Euclidean macro ball.

Finally require the coupling to be small enough that
\begin{equation}
 L_*=\frac{b}{8C_RgT}\ge2,\qquad L_*gT<\rho_R,
 \label{cdm-lambda-radius}
\end{equation}
and the interaction bounds below hold uniformly for
$\epsilon\le\bar\epsilon$. The second condition is an
additional fixed restriction on $b$. For instance $T(g)$
may exceed the fixed threshold and grow as a fixed multiple
of $\sqrt{\log(1/g)}$. Then $\tau(T)\to0$, $gT^3\to0$ and
$gT/b\to0$ for this fixed $b$.

\subsection{Strict coefficients are already in the next background}
\label{cdm-support-section}

Lemma~\ref{cdc-precision-lemma} composes the actual precision
pieces with the supported stationary field while keeping their
matrix rows and columns in the original carrier. Its
scalar-coordinate factor is paid by the independent classical
cardinality reserve. Lemma~\ref{cdc-scalar-lemma} keeps the
whole spectator $\xi$ carrier even at zero background Taylor
degree, joins each substituted background support to it, and
reroots the resulting connected carrier. The classical action
includes its pair-supported quadratic part; its actual chart
Hessian retains both constraint and ordinary-log curvature
terms and the full absolute gauge Hessian.

Thus \eqref{cdc-base-smallness} supplies, directly in $(w,\xi)$,
\begin{equation}
 \norm{V_{\mathrm{base}}}_B\le v,\qquad
 v=C_{\mathrm{base}}(gT^3+gT+g^2T^2).
 \label{cdm-base-bound}
\end{equation}
The independent-background definition subtracts constant,
linear and quadratic Taylor terms before composition. At the
physical branch the total linear term is zero; it is not
silently dropped away from that branch.

For the inherited term, the actual complete-fiber construction
\eqref{cdj-joint-family}--\eqref{cdj-full-family} yields
\begin{equation}
 \norm{f_\lambda-a_p}_{\kappa_0,A_{\mathrm{in}};b,T}
                 \le C_R|\lambda|gT.
 \label{cdm-actual-family}
\end{equation}
Every coefficient reads $\xi$ only in the fixed output-root
collar, independent of its remote $w$ support. On
\eqref{cdm-lambda-radius}, the whole absolute family has
amplitude at most $b/4$. Hence every truncated core has that
same bound without using cancellation.
Theorem~\ref{cdj-transport} and
\eqref{cdj-actual-increment} give canonical connected carriers
\begin{align}
 \sum_K U_{K,\lambda}&=E(f_\lambda),&
 \norm{U_\lambda}_B&\le C_T\epsilon,\notag\\
 C_T&=e^{B5^d2^d+2d}R^dS_1,&
 S_1&=\sum_{z\in\mathbb Z^d}e^{-|z|_1},\notag\\
 \norm{Q_{\mathrm{inh}}}_B
 &\le q:=\frac{C_T\epsilon}{L_*-1}
 \le16C_TC_R\frac{gT}{b}\epsilon .
 \label{cdm-inherited-bound}
\end{align}
The $\lambda=0$ family is termwise $\xi$ independent and
represents $E_{\mathrm{pre}}$. Empty background support terms
of the fluctuation difference remain interactions if they
depend on $\xi$. Banach Cauchy is applied with its contour
integral before the positive anchored sum, as in
\eqref{cdj-Cauchy-increment}.

All these coefficients are already functions of the next
background $w$ on the Gaussian graph. The Gaussian output
will need no second stationary substitution and no second
generic inherited-radius estimate.

\subsection{Nonvanishing and a bound for the inherited logarithm difference}
\label{cdm-nonvanishing-section}

Combine equal-carrier exponents. Their exact multiplier norm is
\begin{equation}
 \sup_z\sum_{K\ni z}e^{B|K|}
       \norm{e^{V_{\mathrm{base},K}+Q_K}-1}
       \le e^{v+q}(v+q).
 \label{cdm-multiplier-bound}
\end{equation}
Products on distinct overlapping carriers are handled by the
overlap theorem. Each multiplier is defined on its actual
carrier window and extended by zero outside it. This is exact
on the global retained window. Since
$B\ge4+\log\mathcal K_\theta+\kappa_G$ and connected $K$
has metric spanning-tree length $|K|-1$, this norm pays the
strict Gaussian multiplier weight.

Impose
\begin{equation}
 v\le\frac{\eta_*}{16e},\qquad
 q\le\frac{\eta_*}{32e},\qquad e^3\tau(T)\le\eta_*/4.
 \label{cdm-quarter-budgets}
\end{equation}
With $\eta_*\le1$, \eqref{cdm-multiplier-bound} is below the
interaction quarter budget. The precision conditions have
already been imposed. The tunable strictly local Gaussian
theorem therefore proves nonvanishing of $Z$ on the full
$b$-ball and an analytic connected logarithm with
\begin{align}
 a_G&=\min\{1/2,\kappa_G/4\}=\frac12,\notag\\
 \norm{\log Z}_{a_G;b}
 &\le2S_1\left[e^3\tau(T)+e^{v+q}(v+q)\right].
 \label{cdm-log-bound}
\end{align}
This is a nonvanishing theorem on a volume-independent
complex domain, not an inference from real positivity.

There is also a precise inherited-log difference estimate.
For $q>0$, replace $Q_{\mathrm{inh}}$ by
$\zeta Q_{\mathrm{inh}}$ and set
\begin{equation}
 L_Q=\frac{\eta_*}{16e q}\ge2.
 \label{cdm-log-interpolation-radius}
\end{equation}
On $|\zeta|<L_Q$, the total exponent norm is at most
$\eta_*/(8e)$; since $\eta_*\le1$, its multiplier norm is
at most $\eta_*/8$. Together with the bad-block budget,
the Gaussian bound is at most $S_1\eta_*$ on that disk.
The canonical connected-log coefficients are holomorphic
as an anchored Banach family. Cauchy therefore gives
\begin{align}
 \norm{\log Z-\log Z_{\mathrm{base}}}_{a_G;b}
 &\le\frac{S_1\eta_*}{L_Q-1}
 \le32eS_1q\notag\\
 &\le512eS_1C_TC_R\frac{gT}{b}\epsilon .
 \label{cdm-log-difference}
\end{align}
For $q=0$ the difference is zero. This interpolation keeps
the base exponent throughout; it does not omit a factor
$e^{V_{\mathrm{base},K}}$ when comparing equal-carrier
multipliers.

\subsection{External normalizations and the exact update}
\label{cdm-update-section}

The pivot and Gaussian determinant corrections are
\begin{equation}
 \Delta(w)=\log\frac{\mathcal D_R(V^\star(w))}
                         {\mathcal D_R(1)},\qquad
 Q_{\det}(w)=-\frac12\log\det[A(w)A_0^{-1}].
 \label{cdm-external-logs}
\end{equation}
The first is the sum of actual local pivot log determinants
composed with the supported stationary field. The scalar
composition lemma gives connected coefficients with any
required fixed reserve and norm at most $C_\Delta b$.

For the second, put $C_0=A_0^{-1}$ and
$\mathcal B(w)=A(w)-A_0$ and expand afresh on the macro graph:
\begin{equation}
 Q_{\det}(w)=\frac12\sum_{n\ge1}\frac{(-1)^n}{n}
                         \operatorname{Tr}[(C_0\mathcal B(w))^n].
 \label{cdm-det-series}
\end{equation}
Join every entire precision label and each covariance path
by deterministic trees and shortest paths. Their cyclic
union is connected and has cardinality at most the total
edge-length budget plus one. The weighted row norm at
$\gamma=8$ and a rerooting spare give
\begin{equation}
 \norm{Q_{\det}}_{1/2;b}
 \le\frac{D_Ge^{1/2}S_1}{2}
       \frac{K_8\nu_8}{1-K_8\nu_8}
 \le C_{\det}b,\qquad K_8\nu_8<1.
 \label{cdm-det-bound}
\end{equation}
Here $K_8$ bounds the macro baseline covariance and $\nu_8$
the strict precision perturbation at that weight, uniformly
in $M$ and volume. Their smallness was included in the
choice of $b$.

With the field-independent factor
\eqref{cdm-normalization}, the exact correction is
\begin{equation}
 E_{\mathrm{next}}=
 E_{\mathrm{pre}}+\Delta+Q_{\det}+\log Z.
 \label{cdm-update}
\end{equation}
Its first term has weight $B$, and the other terms have
weight at least $1/2$, all on the same product ball $b$.
They therefore have the required weight $a_0=1/16$.
Equations \eqref{cdm-inherited-bound},
\eqref{cdm-log-bound} and \eqref{cdm-det-bound} prove
\eqref{cdm-output-bound}. Subtracting each output
coefficient's value at $w=0$ costs at most a factor two
in the norm and changes only the field-independent
normalization.

\subsection{Gauge covariance of the sector and Ward continuation}
\label{cdm-gauge-section}

Take a small real coarse gauge $k$ and lift it to current
sites as a constant on each final $R$-block. A background
with $P_R=1$ keeps that property under this lift. The
composite constraint transforms from $W$ to $W^k$, and
the classical action is invariant on the common
comparison domain. Uniqueness of the nearby physical
constrained minimum thus gives
\begin{equation}
 V^\star(W^k)=
 V^\star(W)^{\,\text{block-constant lift of }k}.
 \label{cdm-background-covariance}
\end{equation}
The relative retained chart transforms by source-endpoint
adjoints: this follows from constraint covariance,
unchanged pivots and uniqueness for the prescribed
retained values.

For real compact gauges these adjoints are orthogonal
in the chosen color norm. They preserve each actual
link/color-ball cutoff and have unit absolute retained
Jacobian. The original Haar/delta formula, equivalently
its exact pulled-back factors, therefore gives equality
of the chosen small-field integrals at $W$ and $W^k$
on a neighborhood of the identity. The inherited
identity is an identity of the total functional with
its compatible branch. No coefficientwise gauge
invariance is used.

The extracted classical action in
\eqref{cdm-background} is also gauge invariant there.
Nonvanishing gives the analytic logarithm on the whole
product ball. Differentiating the real near-origin
identity gives the infinitesimal coarse Ward identity
for $E_{\mathrm{next}}$. Its coefficients are holomorphic
functions of $w$. Uniqueness on the real slice, followed
by the holomorphic identity theorem on the connected
finite-volume ball, extends this identity throughout
$\norm{w}_\infty<b$. Integrating it gives invariance
along gauge paths staying in this ball.

This does not compare arbitrary endpoints whose
connecting gauge path leaves the domain. It also does
not turn the continued germ into a globally identified
current density without the orbit representation
specified above. A field-independent vacuum
normalization leaves the Ward identity unchanged.

\subsection{Finite-step scope}
\label{cdm-scope}

The common radius is obtained by placing the inherited
argument at the actual pre-gauge representative and
expanding only its distant supported tail. All new
classical and measure coefficients are already in
$w$ before integration. The Gaussian output therefore
incurs no second stationary transport. This proves
the stated domain, strict-support and nonvanishing
interfaces of one chosen composite step, uniformly
in preceding homogeneous refinement and volume.

The estimate still contains $C_T\epsilon$, not a
contractive amplitude update. Apart from optional
vacuum normalization, no relevant or marginal
extraction has been performed. A repeated quantum
construction requires a stable extraction and coupling
scheme, amplitude control, large-field and
mixed-history estimates, compatible global measure
identities, and the observable and continuum limits.
None follows here from analyticity or the local
Ward identity alone.

\section{A contraction for the coefficientwise high-field-degree background}
\label{cdh-section}

The following result concerns only the inherited background
$E(a_p(w))$. It uses the actual pre-gauge representative and the
same incoming and outgoing activity exponent and field radius.
It does not identify a gauge-invariant marginal extraction.

Fix $a_0>0$ and a primitive odd $\ell$. Old carriers $H$ are
connected sets of current $J$-cells; output carriers are
$J$-cells of the next current lattice, or current $JR$-cells,
where $R=\ell^s$ and $J=\ell^j$. The projection has fiber $R^d$
and satisfies
\begin{equation}
 |\pi_RH|\leq2^d+\frac{2^d}{k_R}|H|,
 \qquad k_R=1+\left\lfloor\frac{R-3}{4}\right\rfloor.
 \label{cdh-compression}
\end{equation}
Use compatible finite regulators. Write
\begin{equation}
 \|E\|_{a_0,b}
 =\sup_i\sum_{H\ni i}e^{a_0|H|}\|E_H\|_b,
 \label{cdh-input-norm}
\end{equation}
where each coefficient is holomorphic on its scalar-coordinate
radius-$b$ product polydisc. Each old cell has
$D(J)=dJ^d\dim\mathfrak g$ scalar coordinates.

\subsection{Uniform transport on a larger field-scaling disk}
\label{cdh-scaling-section}

Choose $\kappa\geq a_0+1$ and a finite classical reserve
$A\geq d\kappa$ before choosing $R$. The established pre-gauge
theorem gives constants $C_A,\kappa_0,r_A>0$, fixed at primitive
$\ell$ independently of $R,M$ and volume, with
\begin{equation}
 a_{p,i}(w)=\sum_Xa_{i;X}(w_X),\qquad a_{i;X}(0)=0,\qquad
 \sup_i\sum_Xe^{\kappa_0\tau(\{i\}\cup X)+A|X|}
                   \|a_{i;X}\|_r\leq C_A r/R
 \quad(0<r\leq r_A).
 \label{cdh-ap-input}
\end{equation}
The complete current-field maximum in one source $R$-block is
inside the support sum. Put
\begin{equation}
 L_*=\frac{R}{8C_A}\geq2,\qquad
 0<b\leq\frac{r_A}{L_*},\qquad
 f_\lambda(w)=a_p(\lambda w),\quad |\lambda|<L_*.
 \label{cdh-scaling}
\end{equation}
The absolute rooted norm of $f_\lambda$ on the $b$-ball is at
most $b/8$. Every support truncation therefore has that same
component bound without cancellation.

Cut the rooted supports at distance $\ell_{\rm tail}$, using
exponent $\kappa_0/2$ for the tail, and group roots by $J$.
Provided $2\kappa/J\leq\kappa_0/2$, its macro norm obeys
\begin{equation}
 s_{\rm tail}\leq J^d(b/8)e^{-\kappa_0\ell_{\rm tail}/2}.
 \label{cdh-tail}
\end{equation}
The factor $J^d$ is the sum over old output roots before the
macro-fiber maximum. Restriction to an old $J$-current input
fiber introduces no additional coordinate sum.
For $\ell_{\rm tail}\leq J$ use the conservative connected core
$K_0(H)=N_{\infty,2}(\pi_RH)$. Its constants are
\begin{equation}
 c_0=5^d2^d,\qquad c_1=\frac{5^d2^d}{k_R},\qquad
 h_0=2d,\qquad F=R^d,\qquad m=b/2.
 \label{cdh-core-constants}
\end{equation}
Choose $R$ so that $c_1a_0\leq a_0/2$.
Choose $J$ large and
$\ell_{\rm tail}=\lceil c\log J+C\rceil$ with
$c>4d/\kappa_0$, so that $\ell_{\rm tail}\leq J$ and
\begin{equation}
 t:=s_{\rm tail}/m\leq\tfrac12,\qquad
 D(J)\log\frac1{1-t}\leq a_0/2.
 \label{cdh-tail-budget}
\end{equation}
This is possible because
$D(J)t\leq(d\dim\mathfrak g/4)J^{2d}
e^{-\kappa_0\ell_{\rm tail}/2}\to0$.
Larger $J$ required by other fixed-step modules can be used.
These conditions do not depend on $b$.

The joint core/tail theorem, specialized to no fluctuation
variables and target exponent $a_0$, now gives a canonical
connected family $U_\lambda$ with
\begin{equation}
 \sum_KU_{K,\lambda}(w_K)
   =\sum_HE_H(a_p(\lambda w)|_H),\qquad
 \|U_\lambda\|_{a_0,b}\leq C_0R^d\|E\|_{a_0,b},
 \quad
 C_0=e^{a_0 5^d2^d+2d}S_1.
 \label{cdh-uniform-transport}
\end{equation}
Here $S_1=\sum_{z\in\mathbb Z^d}e^{-|z|_1}$ also bounds the
periodic sum. Support choices are fixed independently of
$\lambda$. The family is holomorphic in the anchored activity
Banach space. The prefactor $C_0$ is independent of $J,R,M$ and
volume; $R^d$ is explicit. Although the background substitution
is nonlinear, this canonical construction is linear in $E$.

\subsection{Coefficientwise fifth-order vanishing and contraction}
\label{cdh-contraction-section}

\begin{theorem}[High-field-degree background transport]
\label{cdh-contraction}
Assume, for every input carrier $H$, the full total field
Taylor tensors satisfy
\begin{equation}
 D^nE_H(0)=0,\qquad n=0,1,2,3,4.
 \label{cdh-jet-condition}
\end{equation}
This includes mixed derivatives in all coordinates of $H$.
Under the preceding parameter choices, the connected output
$\mathcal T_{R,J}E:=U_1$ has the same coefficientwise vanishing
jets and obeys
\begin{equation}
 \|\mathcal T_{R,J}E\|_{a_0,b}
 \leq C_0(8C_A)^5R^{d-5}\|E\|_{a_0,b}.
 \label{cdh-general-contraction}
\end{equation}
In $d=4$, for any prescribed $0<q_{\rm contr}<1$, taking the
finite power $R=\ell^s$ sufficiently large that
\begin{equation}
 R\geq C_0(8C_A)^5/q_{\rm contr}
 \label{cdh-R-contraction}
\end{equation}
in addition to the earlier conditions makes this a genuine
$q_{\rm contr}$ contraction. It is uniform in old depth and
compatible volume.
\end{theorem}

\begin{proof}
In the canonical expansion, a term of tail Taylor degree
$q=|\nu|$ contains
$(\partial^\nu E_H)(c_\lambda)/\nu!$ and $q$ supported tail
factors. Core and tail both vanish at $\lambda=0$ termwise.
For $q\leq4$, condition \eqref{cdh-jet-condition} makes the
derivative at the core vanish to order at least $5-q$ in
$\lambda$, and the tail factors supply order $q$.
For $q\geq5$ the tail factors alone supply order at least five.
This includes the core term $q=0$.
Normal convergence on smaller disks permits differentiation
and grouping: every canonical carrier, and hence $U_\lambda$
as a Banach-valued function, has derivatives of orders zero
through four equal to zero. No cancellation between carriers
is used. The same argument in the field variable $w$ proves
the claimed output jet condition.

For completeness, if a Banach-valued holomorphic $F$ on
$|\lambda|<L$ is bounded by $M$ and has a zero of order $p$ at
zero, then
\begin{equation}
 \|F(\lambda)\|\leq M(|\lambda|/L)^p.
 \label{cdh-Banach-Schwarz}
\end{equation}
The singularity of $F(\lambda)/\lambda^p$ is removable.
On a circle of radius $\rho<L$ its norm is bounded by
$M/\rho^p$; apply the scalar maximum principle to every
bounded linear functional and send $\rho\uparrow L$.
Applying this with $p=5$ to
\eqref{cdh-uniform-transport}, at $\lambda=1$, gives
$C_0R^dL_*^{-5}\|E\|$, which is
\eqref{cdh-general-contraction}. A coefficientwise Cauchy
remainder bound would contain an unnecessary
$(1-L_*^{-1})^{-1}$ factor.
\end{proof}

The parameter order is noncircular: choose the finite reserve
and its constants, then $R$, then $J$ and its tail truncation,
and finally the single radius $b\leq8C_A r_A/R$ together with
any other fixed classical restrictions. The same $b$ is used
at the input and output. For the conservative halo in $d=4$
and $a_0=1/16$, $C_0=e^{633}S_1$; no practically sized scale
is asserted. This result concerns the background transport,
not the new quantum correction or the inherited fluctuation
increment.

\subsection{A bounded Taylor subtraction, with its limitations}
\label{cdh-projection-section}

On the circular product polydisc, define the total homogeneous
Taylor projection by
\begin{equation}
 (P_nE_H)(a)=\frac1{2\pi}\int_0^{2\pi}
               e^{-int}E_H(e^{it}a)\,dt.
 \label{cdh-phase-projection}
\end{equation}
It is the total-degree-$n$ Taylor polynomial, has carrier
$H$, and satisfies $\|P_nE_H\|_b\leq\|E_H\|_b$.
The phase integral estimates the complete homogeneous
polynomial without summing its coordinate monomials.
Consequently
\begin{equation}
 \mathcal R_{\geq5}E_H
     =E_H-\sum_{n=0}^4P_nE_H,\qquad
 \|\mathcal R_{\geq5}E\|_{a_0,b}\leq6\|E\|_{a_0,b}.
 \label{cdh-projection-bound}
\end{equation}
No dimension factor, radius shrink or cardinal reserve is
needed for this projection. The remainder meets
\eqref{cdh-jet-condition}, and in $d=4$
\begin{equation}
 \|\mathcal T_{R,J}\mathcal R_{\geq5}E\|_{a_0,b}
 \leq\frac{6C_0(8C_A)^5}{R}\|E\|_{a_0,b}.
 \label{cdh-projected-bound}
\end{equation}

This is not a gauge-invariant extraction scheme. Local gauge
transformations act nonlinearly and affinely in a link-potential
chart, so the low-degree part and high-degree remainder of a
gauge-invariant total need not be separately gauge invariant.
The rooted/pre-gauge identity applies to the complete
inherited functional before such a split; it cannot be applied
separately to these pieces without another argument.

Nor are the low jets automatically a finite-dimensional space
of couplings. They include kernels on arbitrarily many carrier
geometries and separations. Classifying their local
gauge-invariant content, separating low momentum from field
degree, constructing nonlinear Ward-compatible completions,
and bounding marginal or relevant terms remain separate
requirements. New quantum coefficients may regenerate low
jets. The contraction proved here therefore supplies an
irrelevant-background estimate, not a quantum induction,
running coupling or continuum construction.

\section{High-jet differences of the actual quantum step}
\label{cdq-section}

The background contraction can be combined with the actual
common-domain fluctuation integral. The result compares two
admissible complete inherited inputs. It does not assert that
the quantum step preserves their common low jets.

Use the graphs, radius $b$, exponent $a_0=1/16$ and parameters
of Theorem~\ref{cdm-theorem}, with the additional choices below.
Pad absent coefficients by zero so that both inputs have one
common connected-carrier index set. Write
\begin{equation}
 E_j=\sum_HE_{j,H},\qquad
 \norm{E_j}_{a_0;b}\le\bar\epsilon,\qquad
 D=E_1-E_2,\qquad D_H=E_{1,H}-E_{2,H}.
 \label{cdq-inputs}
\end{equation}
Both total inputs satisfy the local gauge and logarithm-branch
hypotheses of \eqref{cdm-input}. Assume the full coefficientwise
condition
\begin{equation}
 D^nD_H(0)=0\qquad(H\text{ arbitrary},\ 0\le n\le4),
 \label{cdq-common-jets}
\end{equation}
including every mixed field derivative. Equality of low jets
only after summing carriers is not sufficient.

\subsection{Compatible coefficient families}
\label{cdq-compatibility-section}

Choose one pre-gauge reserve large enough for the stronger
joint exponent $B=53/8$, and use its same supported
$a_p=\sum_Xa_{p,i;X}$ in Sections~\ref{cdj-section} and
\ref{cdh-section}. Fix one sufficiently large $J$, one
$\ell_{\mathrm{tail}}$, the same support truncation
$X\subset B_{\ell_{\mathrm{tail}}}(i)$, the same halo
$K_0(H)=N_{\infty,2}(\pi_RH)$, and the same deterministic
joining trees and carrier grouping. The stronger joint
budgets imply the weaker background budgets.
Also impose the final-radius condition
\begin{equation}
 L_H=\frac{R}{8C_A}\ge2,\qquad bL_H\le r_A.
 \label{cdq-high-disk}
\end{equation}

In the actual fluctuation family each coefficient has the
form $a_{p,i;X}+d_{\lambda;i;X}$ with
$d_{0;i;X}=0$ termwise. Consequently the core/tail expansion
at fluctuation parameter zero is exactly the background
expansion at field-scaling parameter one:
\begin{equation}
 U^{\mathrm{fluct}}_{K,0}[E]
   =U^{\mathrm{field}}_{K,1}[E]
   =(\mathcal T_{R,J}E)_K.
 \label{cdq-background-identity}
\end{equation}
The local coefficient functions, multiindices, support
occurrences and trees coincide before summing $K$.
The two proofs may estimate this same assignment in different
weights. Extra sites of the halo can be retained as declared
support even when the background coefficient does not read them.

Core/tail differentiation of $E_H$ and support grouping are
linear in the input family. Thus $\mathcal T_{R,J}$ and the
inherited joint difference $Q_E$ are linear, and
\begin{equation}
 Q_{E_1}-Q_{E_2}=Q_D,\qquad
 \norm{Q_E}_B\le a_Q\norm{E}_{a_0;b},\qquad
 a_Q=16C_TC_R\frac{gT}{b}.
 \label{cdq-Q-linear}
\end{equation}
This is the actual estimate \eqref{cdm-inherited-bound}.

For the Gaussian logarithm retain all possible carrier
labels, including coefficients that happen to vanish.
Use common forest and hull assignments, precision pieces
and the normally convergent cluster-log branch. Let
$\mathcal L(Q)$ denote this one analytic coefficient-valued
logarithm map with the fixed base exponent and precision.
The compatible actual output family is
\begin{equation}
 \mathscr F(E)=
 \mathcal T_{R,J}E+\Delta+Q_{\det}+\mathcal L(Q_E).
 \label{cdq-output-map}
\end{equation}
Its scalar sum is precisely \eqref{cdm-update}, with the
same germ or gauge-orbit density interpretation. The
classical background, covariance, cutoff, base factors
and external determinant families are identical for
the two inputs. No $E$-dependent quantum saddle is used.

\subsection{A two-input Lipschitz estimate}
\label{cdq-theorem-section}

Set
\begin{equation}
 C_H=C_0(8C_A)^5,\qquad
 C_{\mathrm{quant}}=1024eS_1C_TC_R,
 \label{cdq-constants}
\end{equation}
where $C_0$ is \eqref{cdh-uniform-transport}. It is fixed
before $R$; $C_{\mathrm{quant}}$ is fixed after the finite
geometry is selected. Strengthen the allowed budgets to
\begin{equation}
 v\le\frac{\eta_*}{64e},\qquad
 a_Q\bar\epsilon\le\frac{\eta_*}{128e},\qquad
 e^3\tau(T)\le\eta_*/4,\qquad \eta_*\le1.
 \label{cdq-budgets}
\end{equation}
All precision and domain conditions of
Theorem~\ref{cdm-theorem} remain in force.

\begin{theorem}[Quantum differences with common low jets]
\label{cdq-theorem}
For the compatible families above, condition
\eqref{cdq-common-jets} implies
\begin{equation}
 \norm{\mathscr F(E_1)-\mathscr F(E_2)}_{a_0;b}
 \le
 \left(\frac{C_H}{R}
             +C_{\mathrm{quant}}\frac{gT}{b}\right)
                       \norm{E_1-E_2}_{a_0;b}.
 \label{cdq-full-bound}
\end{equation}
Writing $\mathcal R_{\ge5}$ for the coefficientwise
total-degree remainder in \eqref{cdh-projection-bound},
\begin{equation}
 \norm{\mathcal R_{\ge5}
       [\mathscr F(E_1)-\mathscr F(E_2)]}_{a_0;b}
 \le6\left(\frac{C_H}{R}
             +C_{\mathrm{quant}}\frac{gT}{b}\right)
                       \norm{E_1-E_2}_{a_0;b}.
 \label{cdq-projected-bound}
\end{equation}
These estimates are uniform in old depth $M$ and compatible
volume. They assert a bound for the specified compatible
canonical families, not for arbitrary localizations of
the same summed functional.
\end{theorem}

\begin{proof}
By \eqref{cdq-background-identity}, the actual background
difference is the family to which
Theorem~\ref{cdh-contraction} applies. The coefficientwise
jet condition gives
\begin{equation}
 \norm{\mathcal T_{R,J}D}_{a_0;b}
                  \le\frac{C_H}{R}\norm{D}_{a_0;b}.
 \label{cdq-background-bound}
\end{equation}
The external $\Delta,Q_{\det}$ cancel coefficientwise.
Linearity in \eqref{cdq-Q-linear} leaves
\[
 \mathscr F(E_1)-\mathscr F(E_2)
   =\mathcal T_{R,J}D
       +\mathcal L(Q_{E_2}+Q_D)-\mathcal L(Q_{E_2}).
\]

If $D=0$ the claim is immediate. Otherwise put
\begin{equation}
 q_D=a_Q\norm{D}_{a_0;b},\qquad
 L_D=\frac{\eta_*}{32e q_D}.
 \label{cdq-interpolation-radius}
\end{equation}
Since $\norm{D}\le2\bar\epsilon$,
\eqref{cdq-budgets} gives $q_D\le\eta_*/(64e)$ and
$L_D\ge2$. Interpolate the inherited exponent as
$Q_{E_2}+\zeta Q_D$ on $|\zeta|<L_D$. Its total exponent
with the fixed base satisfies
\begin{equation}
 \norm{V_{\mathrm{base}}+Q_{E_2}+\zeta Q_D}_B
 \le
 \left(\frac1{64}+\frac1{128}+\frac1{32}\right)
              \frac{\eta_*}{e}
 =\frac{7\eta_*}{128e}<\frac{\eta_*}{8e}.
 \label{cdq-exponent-budget}
\end{equation}
Equal-carrier exponentiation costs at most $e^xx$, where
$x$ is this bound. With $\eta_*\le1$ its multiplier norm
is at most $\eta_*/8$. The tail quarter budget then puts
the connected-log norm at most $S_1\eta_*$ throughout the
disk, by \eqref{cdm-log-bound}. Precision and base remain
fixed, and the cluster theorem supplies one common
nonvanishing logarithm branch.

The interpolated full input $E_2+\zeta D$ need not stay
inside the original $\bar\epsilon$ class on this larger
disk. We have instead verified the actual Gaussian
interaction budgets directly. Normal convergence gives
a holomorphic anchored log family. Banach Cauchy, with
the contour integral before the positive anchored sum,
therefore yields
\begin{align}
 \norm{\mathcal L(Q_{E_1})-\mathcal L(Q_{E_2})}_{a_G;b}
 &\le\frac{S_1\eta_*}{L_D-1}
 \le64eS_1q_D\notag\\
 &\le C_{\mathrm{quant}}\frac{gT}{b}\norm{D}_{a_0;b},
 \qquad a_G=\frac12.
 \label{cdq-log-bound}
\end{align}
Both endpoints are evaluations of the same canonical
analytic cluster family, so no independent logarithm
branch or carrier choice is introduced. Since
$a_G\ge a_0$, this proves \eqref{cdq-full-bound}.
Finally \eqref{cdh-projection-bound} has norm at most
six on the unchanged circular product ball and preserves
the declared carriers. Applying it proves
\eqref{cdq-projected-bound}.
\end{proof}

\subsection{A concrete one-half factor and its limitation}
\label{cdq-scope-section}

The parameter order can give a concrete projected factor.
Fix the stronger common reserve and its $C_A,\kappa_0,r_A$,
then choose a finite power $R=\ell^s$ satisfying all
previous compression and amplitude conditions and
\begin{equation}
 R\ge24C_H.
 \label{cdq-R-choice}
\end{equation}
Choose the common $J$ and tail cutoff afterward, then
one final $b$ satisfying \eqref{cdq-high-disk} and all
classical and Gaussian neighborhoods. Finally choose
$g,T(g)$ to meet \eqref{cdq-budgets} and
\begin{equation}
 C_{\mathrm{quant}}gT/b\le1/24.
 \label{cdq-g-choice}
\end{equation}
With a logarithmically growing cutoff the latter holds
for all sufficiently small $g$ at these fixed parameters.
Equations \eqref{cdq-full-bound}--\eqref{cdq-projected-bound}
then give
\begin{equation}
 \norm{\mathscr F(E_1)-\mathscr F(E_2)}_{a_0;b}
       \le\frac1{12}\norm{D}_{a_0;b},\qquad
 \norm{\mathcal R_{\ge5}
       [\mathscr F(E_1)-\mathscr F(E_2)]}_{a_0;b}
       \le\frac12\norm{D}_{a_0;b}.
 \label{cdq-half-contraction}
\end{equation}
The same $b$ is used at input and output. The canonical
map above uses a common field-independent normalization.
Separate output vacuum subtraction may double the
unprojected estimate, but leaves the projected bound
unchanged because $\mathcal R_{\ge5}$ already kills
constants.

The scalar sums compared here are differences of the
actual common-domain quantum map. The proof first
compares two admissible complete inputs; it does not
apply a rooted/pre-gauge identity separately to a
nongauge-invariant extraction. The representation-sensitive
condition \eqref{cdq-common-jets} and the compatible
family construction are essential to the norm statement.

The full output difference can regenerate field jets
of degrees zero through four through the fluctuation
integral. Thus the quantum map does not preserve the
affine input class with the prescribed low jets.
Projecting the output is bounded but need not preserve
gauge invariance, because gauge transformations are
nonlinear in these potential coordinates. No
gauge-invariant low-jet extraction or closed
relevant/marginal flow has been constructed.
This contractive response is therefore not yet a
Banach fixed-point iteration on admissible physical
activities. Large fields, mixed histories, global
measure consistency, extraction and coupling flow,
and continuum observables remain separate.

\section{The next proof obligation}

The finite step now returns the same analytic radius and the
same positive connected-carrier exponent. Its amplitude estimate
still has a coefficient larger than one multiplying the inherited
activity norm. The background-dependent determinant and delta
terms are also not small merely because the coupling is small.
These statements locate the remaining small-field induction
problem: extract the vacuum and allowed marginal terms, establish
the coupling recursion, and prove a uniform bound on the
remaining interactions through arbitrarily many steps.

The high-field-degree estimates provide one part of that
recursion: coefficientwise fifth-order background remainders
contract at a sufficiently large fixed macro ratio, and
the actual quantum step has a contractive response to
differences with matching low jets. Its Taylor subtraction
is bounded without a radius loss. This does not give an
invariant high-degree input class, because the fluctuation
integral can regenerate lower-degree terms.
The low-degree pieces still include kernels on arbitrarily
many carrier geometries, and the separate pieces need not be
gauge invariant. A complete extraction must combine the
field-degree estimate with spatial locality and the nonlinear
gauge identities, and include the generated low jets in the
recurrence.

The total correction has the local gauge identity needed for
that analysis. A scalar marginal gauge-action coefficient still
requires the appropriate lattice symmetry classification; a
Ward identity alone does not supply it. The large-field
complements, mixed-history matching, observable estimates,
continuum nontriviality and reconstruction remain separate
required parts of the four-dimensional existence theorem.

\begin{thebibliography}{9}
\bibitem{priorcd}
\emph{Toward a four-dimensional Yang--Mills construction},
research Notebooks 15--27 in this workspace.
The homogeneous classical support families, strict overlapping
Gaussian expansion, gauge identities, supported pre-gauge
representative, auxiliary grouping and exact composite macrostep
are the principal inputs.
\bibitem{balaban116cd}
T.~Ba\l aban,
\emph{Renormalization group approach to lattice gauge field
theories. II. Cluster expansions},
Commun. Math. Phys. \textbf{116} (1988), 1--22.
\href{https://scispace.com/pdf/renormalization-group-approach-to-lattice-gauge-field-4penm4k7hc.pdf}{Original article copy}.
\bibitem{sourcecd}
E.~Sussman and collaborators,
\href{https://github.com/ethanwsussman/Yang--Mills-LLM-experiment}{Yang--Mills-LLM-experiment},
deposited raw accounts at commit
\texttt{e392ef4501\allowbreak 3a581664ba\allowbreak b76b908671\allowbreak 70b30fe10e}.
\end{thebibliography}
\end{document}
